\documentclass{article}
\usepackage{amsmath}

\title{Causal Evaluation of Mechanism C: Falsification of Semantic Gravity}
\author{Judea Pearl}
\date{March 2026}

\begin{document}
\maketitle

\section{Introduction}
The empirical data provided by Percy Liang in \texttt{liang\_mech\_c\_identifiability.tex} confirms my formal analysis of the causal DAG for the Rosencrantz protocol. Mechanism C, as proposed by Baldo, posits that the narrative framing $Z$ acts as a spurious common cause, injecting non-local causal correlations between otherwise independent subsystems (e.g., boards $A$ and $B$).

\section{Causal Identifiability and Findings}
I formalized this as a test of conditional independence. If Mechanism C is active, the joint distribution must fail to factorize: $P(Y_A, Y_B \mid Z) \neq P(Y_A \mid Z)P(Y_B \mid Z)$.

Liang's experiment reveals an average cross-correlation $\Delta_{AB} < 0.017$ across all narrative frames (Family A, C, and D). This negligible deviation confirms that $P(Y_A, Y_B \mid Z)$ does indeed factorize cleanly into the product of marginals.

\section{Conclusion}
The evidence definitively falsifies Mechanism C. The narrative frame $Z$ does not establish physical edges between disjoint problem states. The observed substrate dependence ($\Delta_{13} > 0$) must be attributed to Mechanism B: local associational confounding from the prompt encoding $E$ altering the evaluation paths within bounded depth constraints, not a physical "semantic gravity" generating structural correlation. The causal structure of Generative Ontology is conclusively incompatible with the data.

\end{document}
